%% s4_1_detail.tex
%% M78 sub-agent output, 2026-05-06
%% §4.1 detail: Wedge algebra A_obs is type-III_1 (TBD-1 sketch)
%% Ready to merge into bianchi_ix_modular_shadow.tex at \subsection{Step A}
%%
%% Refs used:
%%  [BW76]     Bisognano-Wichmann (1976), Commun. Math. Phys. 45:199--224 (no arXiv)
%%  [CLPW2022] arXiv:2206.10780 (Chandrasekaran-Longo-Penington-Witten), WebFetch-verified 2026-05-06
%%  [KLS2024]  arXiv:2406.01669 (Kudler-Flam-Leutheusser-Satishchandran), WebFetch-verified 2026-05-06
%%  [Witten2022] arXiv:2112.12828, WebFetch-verified 2026-05-06
%%  [Speranza2025] arXiv:2504.07630, WebFetch-verified 2026-05-06
%%  [HartnollDCS2024] arXiv:2312.11622, WebFetch-verified 2026-05-06
%%
%% Note: KLS 2024 result (FLRW algebra directly type-II_infty via KMS, not III_1-then-crossed)
%% is accurately reflected below; the III_1 step for BIX is therefore a GENUINE open problem
%% and not a routine analogy.

% ============================================================
\subsection{Step~A: Wedge algebra and its type (TBD-1 detail)}
\label{ssec:wedge-detail}
% ============================================================

\subsubsection*{Setup}

Let $M_{\mathrm{BIX}} = \mathbb{R}\times S^3$ carry a Bianchi~IX vacuum metric
\begin{equation}
  ds^2 = -N^2(\tau)\,d\tau^2
         + \tfrac{1}{4}\,e^{-2\Omega(\tau)}\,
           \bigl(e^{2\beta(\tau)}\bigr)_{ij}\,\sigma^i\otimes\sigma^j,
  \label{eq:BIX-metric}
\end{equation}
where $\tau$ is Misner volume time ($\Omega = -\frac{1}{3}\ln V$,
monotone in the BKL backward direction),
$\{\sigma^i\}$ are the left-invariant 1-forms on $S^3\cong\mathrm{SU}(2)$,
and $\beta = \beta_+\,\mathrm{diag}(1,1,-2) + \beta_-\,\mathrm{diag}(1,-1,0)$
is the Misner anisotropy matrix.
The spatial sections are homogeneous under the SO(3) left-action.
In the BKL (Mixmaster) phase approaching the singularity
$\tau\to+\infty$, the metric~\eqref{eq:BIX-metric} undergoes an infinite
sequence of Kasner epochs separated by bounces encoded by the
Gauss-shift map $\sigma_G(x)=\{1/x\}$~\cite{ManinMarcolli2015}.

Fix an intrinsic observer $\gamma:\mathbb{R}\to M_{\mathrm{BIX}}$
with worldline parametrised by $\tau$.
Define the \emph{wedge complement}
\begin{equation}
  W_\gamma = \{\,p\in M_{\mathrm{BIX}}\;:\;
     p \text{ is spacelike-separated from every point on }\gamma\,\}
\end{equation}
(well-defined locally, since the homogeneous SO(3) foliation provides a
canonical spacelike splitting~\cite{HeinzleUggla2009}).
Consider a free real scalar field $\phi$ on $M_{\mathrm{BIX}}$
satisfying the covariant Klein--Gordon equation $\Box_g\phi = m^2\phi$;
the graviton case is analogous but algebraically more involved.
Let $\mathcal{A}(W_\gamma)$ be the von~Neumann algebra of bounded
operators generated by the Weyl operators $W(f) = e^{i\phi(f)}$
with $\mathrm{supp}(f)\subset W_\gamma$, in the GNS representation
of a reference state~$\omega_{\mathrm{BIX}}$ (specified below).

\begin{conjecture}[Wedge algebra type, M45.1.A --- TBD-1]
\label{conj:wedge-type}
$\mathcal{A}(W_\gamma)$ is the unique hyperfinite type-\textrm{III}$_1$ factor
$\mathcal{R}_\infty$.
\end{conjecture}

\paragraph{Why the standard arguments do not apply.}
The classical route to type-III$_1$ for wedge algebras
runs through the \emph{Bisognano--Wichmann theorem}~\cite{BW76}:
in Minkowski space, the modular Hamiltonian of the Rindler wedge in the
Minkowski vacuum is the \emph{boost generator}~$K$, a global Killing
vector field.
By Tomita--Takesaki theory, this forces the modular automorphism
group $\sigma_t^{\omega}(A) = e^{itK}Ae^{-itK}$ to be a continuous
one-parameter group of outer automorphisms, and the associated Connes
spectrum $S(\mathcal{A}) = [0,\infty)$, which characterises
type-III$_1$~\cite{Witten2022}.
For the de~Sitter static patch,
Chandrasekaran--Longo--Penington--Witten~\cite{CLPW2022}
show the algebra is the unique hyperfinite III$_1$ factor
$\mathcal{R}_\infty$, again by identifying the
modular Hamiltonian with the de~Sitter isometry generator
(a Killing vector of the static patch).

\emph{Neither argument transfers to Bianchi~IX.}
The metric~\eqref{eq:BIX-metric} admits \emph{no global timelike Killing
vector}: the three spatial Killing vectors $\{\xi_i\}$ generating the
SO(3) symmetry are spacelike, and the Misner volume~$\tau$ is not a
Killing coordinate (it appears only as a synchronous gauge choice).
Consequently, the modular Hamiltonian of $\mathcal{A}(W_\gamma)$
is \emph{not} identified with any known geometric symmetry generator,
and the Bisognano--Wichmann route is unavailable.

We note also that the FLRW case handled by
Kudler-Flam--Leutheusser--Satishchandran~\cite{KLS2024}
follows a \emph{different} path: the FLRW algebra is shown to be
type-II$_\infty$ \emph{directly} (via the inflaton zero-mode and
the KMS condition with respect to conformal-time translations),
not via a prior type-III$_1$ identification.
The Bianchi~IX situation is harder: the anisotropic $\beta$-dynamics
breaks the FLRW conformal-time symmetry, so the KLS shortcut is also
unavailable.

\paragraph{Three partial strategies toward TBD-1.}

\vspace{4pt}
\noindent\textbf{Strategy~(i): Haag--Kastler axioms $+$ SO(3) homogeneity.}
If the net $O\mapsto\mathcal{A}(O)$ satisfies the Haag--Kastler
axioms (isotony, locality, covariance under the residual SO(3) isometry)
and the Bianchi~IX quantum field theory satisfies the
\emph{Buchholz--Junglas split property} for the homogeneous spatial foliation
(i.e., for any two causally complete regions $O_1\Subset O_2$,
there exists an intermediate type-I factor),
then standard AQFT theorems~\cite{Buchholz1986,BV1996} imply
that the local algebras are type-III factors.
The additional step from type-III to type-III$_1$ requires showing
that the Connes $S$-invariant $S(\mathcal{A}(W_\gamma))=[0,\infty)$;
we argue this below via the BKL ergodicity.

\noindent\textbf{[TBD-1.1: prove]} The split property for free
scalar fields on Bianchi~IX has not been established.
The proof would require uniform energy bounds on the quantum field
relative to the BIX homogeneous foliation ---
a non-trivial task given the time-dependent metric~\eqref{eq:BIX-metric}
and the absence of a static time.

\vspace{4pt}
\noindent\textbf{Strategy~(ii): Connes $S$-invariant via BKL ergodicity.}
For any factor $\mathcal{N}$ with faithful normal weight $\phi$,
the Connes $S$-invariant is defined as
\begin{equation}
  S(\mathcal{N}) \;=\; \bigcap_{\phi}\,
  \mathrm{Sp}(\Delta_\phi),
  \label{eq:S-invariant}
\end{equation}
where $\Delta_\phi$ is the modular operator of $\phi$ and
the intersection runs over all faithful normal semifinite weights.
$\mathcal{N}$ is type-III$_1$ if and only if $S(\mathcal{N})=[0,\infty)$,
type-III$_\lambda$ (for $\lambda\in(0,1)$) if $S(\mathcal{N})=
\{\lambda^n:n\in\mathbb{Z}\}$, and type-III$_0$ if $S(\mathcal{N})=\{0,1\}$.

For the reference state $\omega_{\mathrm{BIX}}$,
the modular operator $\Delta_{\omega_{\mathrm{BIX}}}$ encodes the
``temperature'' seen by the observer along $\gamma$.
The Kasner-bounce sequence in the BKL phase is encoded by the
Gauss-shift $\sigma_G$ with Kolmogorov--Sinai entropy
$\lambda_{\mathrm{BKL}} = \pi^2/(6\log 2) > 0$
(Definition~\ref{def:lambda-BKL}).
The positive KS entropy means the Gauss-shift is
\emph{strongly mixing} and \emph{Bernoulli} (Adler--Weiss theorem),
so the associated time-evolution has \emph{absolutely continuous
spectrum} on $L^2([0,1],\mu_{\mathrm{Gauss}})$.
If this continuous spectrum is reflected in the modular spectrum
$\mathrm{Sp}(\Delta_{\omega_{\mathrm{BIX}}})$
via the Manin--Marcolli geodesic-flow representation~\cite{ManinMarcolli2015},
one expects $S(\mathcal{A}(W_\gamma))=[0,\infty)$, i.e., type-III$_1$.

\noindent\textbf{[TBD-1.2: prove]} The link between BKL Gauss-shift KS entropy
(a classical dynamical systems result) and the quantum modular spectrum
$\mathrm{Sp}(\Delta_{\omega_{\mathrm{BIX}}})$ is conjectural.
Making it precise requires:
\begin{enumerate}
  \item[(a)] Identifying $\omega_{\mathrm{BIX}}$ as a KMS state
     (or quasi-free state) for some BIX-adapted dynamics;
  \item[(b)] Showing the modular operator $\Delta_{\omega_{\mathrm{BIX}}}$
     has continuous spectrum via the Manin--Marcolli embedding
     $\mathcal{A}(W_\gamma)\hookrightarrow
     B(L^2(X(2),\mu_{\mathrm{Gauss}}))$
     (which is itself Conjecture~M45.1.C);
  \item[(c)] Ruling out III$_\lambda$ ($\lambda\in(0,1)$) ---
     i.e., that the spectrum is the full $[0,\infty)$ and not
     a discrete subgroup $\{\lambda^n\}$.
\end{enumerate}
For step~(c), the aperiodicity of the Gauss shift (its spectrum is purely
continuous on $L^2$ with no point masses at $\{\lambda^n\}$) is encouraging
but does not automatically imply III$_1$ for the quantum algebra.

\vspace{4pt}
\noindent\textbf{Strategy~(iii): Mixmaster scrambling as surrogate for geometric modular flow.}
In de~Sitter, the de~Sitter temperature $T_{dS} = H/2\pi$ controls the
modular flow rate, with $T_{dS}$ arising from the Killing horizon structure.
For Bianchi~IX, there is no Killing horizon, but the Mixmaster oscillations
produce an effective ``scrambling temperature'' at rate $\lambda_{\mathrm{BKL}}$.
One may \emph{conjecture} (heuristically) that the modular flow of
$\mathcal{A}(W_\gamma)$ in $\omega_{\mathrm{BIX}}$ is
\emph{ergodic relative to the Gauss-measure foliation of $X(2)$}, and
that this ergodicity --- in the quantum context --- forces
$S(\mathcal{A}(W_\gamma)) = [0,\infty)$.

\noindent\textbf{[TBD-1.3: prove]} This heuristic requires:
identifying a \emph{geometric} realisation of the modular Hamiltonian
$K_{\mathrm{BIX}}$ as a differential operator on
$L^2(X(2),\mu_{\mathrm{Gauss}})$;
and showing $K_{\mathrm{BIX}}$ has purely continuous spectrum.
A possible route is the Bost--Connes arithmetic algebra associated with
$\Gamma(2)$ (Connes--Marcolli, \emph{Noncommutative Geometry, Quantum
Fields and Motives}, Ch.~III), where the modular Hamiltonian is related
to the geodesic flow generator on $X(2)$ and the spectrum is known to be
continuous.
This route depends on Conjecture~M45.1.C and is therefore circular
without an independent proof of Step~C.

\subsubsection*{Summary of TBD-1 sub-markers}

\begin{center}
\small
\begin{tabular}{lll}
\toprule
Marker & Dependency & Status \\
\midrule
TBD-1.1 & Split property for free scalar on BIX homogeneous foliation
        & Open (AQFT/curved spacetime) \\
TBD-1.2 & Gauss-shift continuous spectrum $\Rightarrow$ $S(\mathcal{A})=[0,\infty)$
        & Open (requires steps (a)--(c) above) \\
TBD-1.3 & BIX modular Hamiltonian = geodesic flow generator on $X(2)$
        & Open (circular with Conj.~M45.1.C) \\
\bottomrule
\end{tabular}
\end{center}

\noindent The most tractable near-term route is
Strategy~(i) $+$ partial~(ii):
establish the Haag--Kastler net and split property on BIX homogeneous
slices (using SO(3) covariance and energy bounds for the free scalar
in a time-dependent homogeneous background), obtain the type-III label,
and then use the Gauss-shift ergodicity as heuristic input to argue
III$_1$ rather than III$_\lambda$.
A rigorous proof of type-III$_1$ for any non-Killing-symmetric cosmological
background would be a significant result in AQFT on curved spacetime
in its own right, independent of the applications to crossed products
pursued here.
